\documentclass{situationist-preprint}

\title{The Scope of Hume's Law}
\msid{MTE-10-b}
\msfamily{MTE-10}
\msdiscipline{Metaethics}
\mstype{Companion manuscript}
\mskeywords{Hume's law; is and ought; normative bridge premise; Searle; thick concepts; naturalism; barriers to implication}

\begin{document}
\maketitle

\begin{abstract}
MTE-10-a proves Hume's law in a corrected form: no conclusion that is normatively significant given its premises follows from purely descriptive premises, so that ought follows from is only with a normative bridge premise. This note fixes the scope of the corrected law. We show that it is a thesis about logical consequence and not about metaphysics, so that it neither supports nor threatens naturalism; that the a posteriori identity of reasons with natural facts is itself a bridge premise; that Searle's derivation of ought from promising, and derivations through thick concepts, succeed only by containing bridge premises; and that the law has a mirror image, on which the principle that ought implies can is a bridge principle and not a law of logic.
\end{abstract}

\section{The corrected law}

MTE-10-a works with admissible models $(d,n)$ that pair a descriptive state with a normative state. Its one substantive assumption is logical autonomy: if $(d,n)$ and $(d',n')$ are admissible, so is $(d,n')$ (Assumption~2.1 of MTE-10-a). A sentence is \emph{purely descriptive} if admissible models with the same descriptive state agree on it (Definition~2.2); a conclusion is \emph{normatively significant} given premises if some model of the premises can be varied in its normative state alone so as to falsify the conclusion, and a \emph{bridge premise} is a premise that is not purely descriptive (Definition~2.3). Theorem~3.1 of MTE-10-a is that purely descriptive premises entail no conclusion normatively significant given them, and Corollary~3.2 that any argument whose conclusion is normatively significant contains a bridge premise.

The classification is semantic, not syntactic. A sentence counts as a bridge premise because of how its truth depends on the normative state, not because it contains a normative word. This is what makes the law immune to Prior's counterexamples \citep{Prior1960}, and it is also what determines its scope. We take four questions of scope in turn.

\section{Logic, not metaphysics}

Logical autonomy concerns admissible models, which represent what logic leaves open. It does not say that every combination of descriptive and normative facts is metaphysically possible. The naturalist of MTE-01 holds that moral facts are natural facts about idealized reactions, so that, metaphysically, the descriptive state fixes the normative state; and this is consistent with logical autonomy, since the dependence is not a matter of logic. The corrected law is therefore neutral on the metaphysical question that Moore thought it settled \citep{Moore1903}. A metaphysical supervenience of the normative on the descriptive no more violates the law than the metaphysical necessity of \enquote{water is H$_2$O} makes the inference from \enquote{this is H$_2$O} to \enquote{this is water} logically valid.

\begin{proposition}\label{p:identity}
The a posteriori identity of reasons with natural facts, Principle~P4 of MTE-11-a, is a normative bridge premise.
\end{proposition}

\begin{proof}
Principle~P4 states that for a fact to be a reason for an agent to act is for it to be such that idealized reactions would favor the act in response to it. Take an admissible model $(d,n)$ at which the identity holds and some fact is a reason; by logical autonomy, $(d,n')$ is admissible for a normative state $n'$ in which nothing is a reason. The idealized reactions are part of the descriptive state $d$ and are unchanged, so at $(d,n')$ the identity fails. Two models with the same descriptive state disagree on it, so it is not purely descriptive.
\end{proof}

The proposition explains how the naturalist derives moral conclusions from natural facts without violating Hume's law. From \enquote{the idealized reactions would disfavor this act} alone, nothing normatively significant follows. Add the identity, and the conclusion that there is a reason not to perform the act follows, as Corollary~3.2 of MTE-10-a permits. The identity is known a posteriori, if at all, and it is the naturalist's bridge. This is also why the identity is not a definition of reasons in other normative terms, as Proposition~6.1 of MTE-11-a requires: it links the normative to the descriptive, not one normative notion to another.

\section{Institutional and thick premises}

\paragraph{Promising.} \citet{Searle1964} derives \enquote{Jones ought to pay Smith five dollars} from \enquote{Jones uttered the words \enquote{I hereby promise to pay you, Smith, five dollars}}, by way of \enquote{Jones promised}, \enquote{Jones placed himself under an obligation} and \enquote{Jones is under an obligation}. The derivation is valid with the premises Searle supplies, and its conclusion is normatively significant. By Corollary~3.2 of MTE-10-a it must contain a bridge premise, and the corollary locates it. Either \enquote{Jones promised} is read as purely descriptive, recording that a certain speech act was performed under the conventions of promising, in which case the step to \enquote{Jones placed himself under an obligation} requires the premise that one who promises is obliged, whose truth varies with the normative state; or \enquote{Jones promised} is read as already including that Jones undertook an obligation he is under, in which case it is itself not purely descriptive. Searle's observation that the obligation is constituted by an institution is correct, but the claim that the institution's obligations bind is not settled by the descriptive facts about the institution, as Hare objected \citep{Hare1964}.

\paragraph{Thick concepts.} An argument from \enquote{the remark was lewd} to \enquote{there was reason not to make it} has a normatively significant conclusion. On the separabilist view, \enquote{lewd} factors into a descriptive and an evaluative component, and the premise is a conjunction of which one conjunct is a bridge premise. MTE-15 argues that thick concepts do not factor. The corrected law does not need them to: it classifies sentences by their truth conditions, and whether or not \enquote{lewd} factors, the truth of \enquote{the remark was lewd} varies with the normative state across models with the same descriptive state, since the remark could have been sexual in every descriptive respect while lacking the evaluative feature. Sentences containing thick terms are therefore bridge premises, whether or not their meaning decomposes. The inseparability of thick concepts makes the bridge harder to see, not absent.

\section{The mirror image}

Logical autonomy is symmetric. If $(d,n)$ and $(d',n')$ are admissible, then so are $(d,n')$ and, exchanging the roles of the two models, $(d',n)$. Call a sentence \emph{purely normative} if admissible models with the same normative state agree on it, and a conclusion \emph{descriptively significant} given premises if some model of the premises can be varied in its descriptive state alone so as to falsify it. The proof of Theorem~3.1 of MTE-10-a, with the roles of $d$ and $n$ exchanged, shows that purely normative premises entail no conclusion descriptively significant given them. In particular, \enquote{Jones ought to pay Smith} does not by itself entail \enquote{Jones can pay Smith}: the principle that ought implies can, whatever its standing, is a bridge principle and not a law of logic. The corrected law and its mirror image are instances of a general pattern of barriers to implication, which \citet{RussellRestall2010} establish for the universal and the particular, the present and the future, and the necessary and the contingent as well as for is and ought. In each case the barrier holds for premise sets that lie wholly on one side of a division whose two sides logic allows to vary independently.

\section{Objections and replies}

\begin{objection}
If even the naturalist's identity is a bridge premise, then the law is trivial: it says only that one cannot derive the normative from the descriptive without a premise linking them.
\end{objection}

\begin{reply}
It says more. The usual formulation is false, as Prior showed, and the corrected law locates the needed link in the premises' truth conditions rather than in their vocabulary. It thereby rules out three claims that are often made: that the absence of normative words from the premises guarantees the failure of the inference, which Prior's second argument refutes; that a premise free of normative words is purely descriptive, which Searle's institutional premise refutes; and that a derivation from natural facts refutes naturalism or establishes it, since the identity premise is needed and is neither logical nor normatively idle.
\end{reply}

\section{Conclusion}

Hume's law in the corrected form of MTE-10-a is a thesis about logical consequence under logical autonomy. It holds for purely descriptive premise sets and applies to sentences according to their truth conditions. Within that scope, Searle's derivation, derivations through thick concepts and the naturalist's derivation of reasons from idealized reactions all succeed only because they contain bridge premises; the last of these is the a posteriori identity of reasons with natural facts. Outside that scope, the law says nothing about metaphysics, and nothing against naturalism.

\begin{thebibliography}{9}\small
\bibitem[Hare(1964)]{Hare1964} Hare, R.~M. (1964). The promising game. \emph{Revue Internationale de Philosophie}, 18(70), 398--412.
\bibitem[Moore(1903)]{Moore1903} Moore, G.~E. (1903). \emph{Principia Ethica}. Cambridge University Press.
\bibitem[Prior(1960)]{Prior1960} Prior, A.~N. (1960). The autonomy of ethics. \emph{Australasian Journal of Philosophy}, 38(3), 199--206.
\bibitem[Russell and Restall(2010)]{RussellRestall2010} Russell, G. and Restall, G. (2010). Barriers to implication. In C. Pigden (ed.), \emph{Hume on Is and Ought}, 243--259. Palgrave Macmillan.
\bibitem[Searle(1964)]{Searle1964} Searle, J.~R. (1964). How to derive \enquote{ought} from \enquote{is}. \emph{The Philosophical Review}, 73(1), 43--58.
\end{thebibliography}

\end{document}
